\section{Classes of Subsets}

\noindent
Let $\Omega$ be a set and let $\Po(\Omega)$ denote the set of all subsets of $\Omega$. 

\begin{definition}
  Let $\Sa$ be a subset of $\Po(\Omega)$. $\Sa$ is a \textbf{semi-algebra} if:
  \begin{enumerate}
    \item $\Omega \in \Sa$;
    \item if $A, B \in \Sa$, then $A \cap B \in \Sa$; and,
    \item if $A \in \Sa$, then there exist some $A_1, A_2, ..., A_n \in \Sa$, where $A_i \cap A_j = \varnothing$ if $i \neq j$ and $A^c = \bigsqcup_{i = 1}^{n} A_i$.
  \end{enumerate}
\end{definition}

\begin{example}
  Let $\Omega = \R$ and $\Sa = \{ \varnothing, \R, (a, b], (-\infty, b], (a, \infty) : a, b \in \R \text{ and } a < b \}$. $\Sa$ is a semi-algebra.
\end{example}

\begin{definition}
  Let $\A$ be a subset of $\Po(\Omega)$. $\A$ is an \textbf{algebra} if:
  \begin{enumerate}
    \item $\Omega \in \A$;
    \item if $A, B \in \A$, then $A \cap B \in \A$; and,
    \item if $A \in \A$, then $A^c \in \A$.
  \end{enumerate}
\end{definition}

Note that this definition implies that an algebra is also closed under finite unions. That is, if $A, B \in \A$, then $A \cup B \in \A$; for $A, B \in \A$ implies $A^c, B^c \in \A$, and by closure under intersection we have $A^c \cap B^c \in \A$, i.e., $(A \cup B)^c \in \A$. Hence, $A \cup B \in \A$.

\begin{definition}
  Let $\F$ be a subset of $\Po(\Omega)$. $\F$ is a \textbf{$\sigma$-algebra} if:
  \begin{enumerate}
    \item $\Omega \in \A$;
    \item if $A_j \in \F$ for all $j \geq 1$, then $\bigcup_{j \geq 1} A_j \in \F$; and,
    \item if $A \in \F$, then $A^c \in \F$.
  \end{enumerate}
\end{definition}

Again, this definition of a $\sigma$-algebra implies closure under countable intersections as well.\\

\noindent Observe that:
\begin{itemize}
  \item Every algebra is a semi-algebra.
  \item Every $\sigma$-algebra is an algebra.
\end{itemize}

\begin{remark}
  \begin{enumerate}
    \item Let $\A_{\alpha} \sub \Po(\Omega)$, $\alpha \in I$, be a family of algebras, where $I$ is an index set, possibly uncountable. Then, $\bigcap_{\alpha \in I} \A_{\alpha}$ is also an algebra.

    \item Similarly, if $\F_{\alpha} \sub \Po(\Omega)$, $\alpha \in I$ is a family of $\sigma$-algebras, where $I$ is a possibly uncountable index set, then, $\bigcap_{\alpha \in I} \F_{\alpha}$ is also a $\sigma$-algebra.
  \end{enumerate}
\end{remark}

\begin{definition}
  Let $\C \sub \Po(\Omega)$. Then, an algebra $\A(\C)$ is the \textbf{algebra generated by $\C$} if $\C \sub \A(\C)$ and $\A(\C) \sub \mathscr{B}$ for any algebra $\mathscr{B}$ containing $\C$ (i.e., $\C \sub \mathscr{B}$).
\end{definition}

$\A(\C)$ is the ``smallest'' algebra containing $\C$. Such an algebra exists, for $\Po(\Omega)$ is an algebra which contains $\C$, and by Remark 2.5, the intersection of all algebras containing $\C$ is also an algebra. This is precisely the algebra generated by $\C$. 

\begin{definition}
  Let $\C \sub \Po(\Omega)$. Then, a $\sigma$-algebra $\sigma(\C)$ is the \textbf{$\sigma$-algebra generated by $\C$} if $\C \sub \sigma(\C)$ and $\sigma(\C) \sub \F$ for any $\sigma$-algebra $\F$ containing $\C$ (i.e., $\C \sub \F$).
\end{definition}

Again, $\sigma(\C)$ exists, for $\Po(\Omega)$ is a $\sigma$-algebra containing $\C$ and by Remark 2.5, the intersection of all $\sigma$-algebras containing $\C$ is also a $\sigma$-algebra, which is in fact the $\sigma$-algebra generated by $\C$. 

\begin{lemma}
  Let $\Sa \sub \Po(\Omega)$ be a semi-algebra. The algebra generated by $\Sa$, $\A(\Sa)$, is the collection of all subsets of $\Omega$ that can be expressed as a finite disjoint union of sets in $\Sa$. That is, $A \in \A(\Sa)$ if and only if there exist disjoint $E_j \in \Sa$, $1 \leq j \leq n$, such that $A = \bigsqcup_{j = 1}^n E_j$. 
\end{lemma}

\begin{proof}
  We first show that if $A = \bigsqcup_{j = 1}^n E_j$ for some disjoint $E_j \in \Sa$, then $A \in \A(\Sa)$. Since $E_j \in \Sa$ and $\Sa \sub \A(\Sa)$, $E_j \in \A(\Sa)$. $\A(\Sa)$ is an algebra and so $\bigsqcup_{j = 1}^n E_j \in \A(\Sa)$.

  We now show the reverse direction. \textcolor{red}{Define
  \[
  \mathscr{B} = \left\{ \bigsqcup_{j = 1}^n F_j : F_j \in \Sa \right\}.
  \]
  }

  Then, $\mathscr{B}$ is the collection of all subsets of $\Omega$ which are finite disjoint unions of sets in $\Sa$. \textcolor{red}{We show that $\mathscr{B}$ is an algebra.}

  Clearly, $\Sa \sub \mathscr{B}$. Since $\Omega \in \Sa$, $\Omega \in \mathscr{B}$.

  \textbf{Closure under finite intersections.} Let $A, B \in \mathscr{B}$. Then, $A = \bigsqcup_{j = 1}^n E_j$ and $B = \bigsqcup_{k = 1}^m F_k$ for some $E_j, F_k \in \Sa$. Now, 
  \[
  A \cap B = \left( \bigsqcup_{j = 1}^n E_j \right) \cap \left( \bigsqcup_{k = 1}^m F_k \right) = \bigsqcup_{j = 1}^n \bigsqcup_{k = 1}^m (E_j \cap F_k)
  \]
  where the second equality is due to the distributive property: $\left( \bigsqcup_{j = 1}^n E_j \right) \cap \left( \bigsqcup_{k = 1}^m F_k \right)$ is the disjoint union over all pairs $(E_j \cap F_k)$. It is a disjoint union because $(E_{j_1} \cap F_{k_1}) \cap (E_{j_2} \cap F_{k_2}) = \varnothing$, as $E_{j_1}, E_{j_2}$ are disjoint and $F_{k_1}, F_{k_2}$ are disjoint. Hence, $A \cap B$ is a finite disjoint union of sets in $\Sa$ and so $A \cap B \in \mathscr{B}$, i.e., $\mathscr{B}$ is closed under finite intersections. 

  \textbf{Closure under complementation.} Now suppose $A \in \mathscr{B}$. Then, $A = \bigsqcup_{j = 1}^n E_j$ where $E_j \in \Sa$. Therefore, 
  \[
  A^c = \left( \bigsqcup_{j = 1}^n E_j \right)^c = \bigcap_{j = 1}^n E_j^c = E_1^c \cap ... \cap E_n^c.
  \]
  Since $E_j \in \Sa$, $E_j^c = \bigsqcup_{k_i = 1}^{l_i} F_{i, k_i}$ for some $F_{i, j} \in \Sa$. Hence, 
  \begin{align*}
  A^c &= \left( \bigsqcup_{k_1 = 1}^{l_1} F_{1, k_1} \right) \cap ... \cap \left( \bigsqcup_{k_n = 1}^{l_n} F_{n, k_n} \right) \\
  &= \bigsqcup_{k_1 = 1}^{l_1} ... \bigsqcup_{k_n = 1}^{l_n} (F_{1, k_1} \cap ... \cap F_{n, k_n})
  \end{align*}
  by the distributive property. Again, $F_{1, k_1} \cap ... \cap F_{n, k_n} \in \Sa$ and they are pairwise disjoint. Thus, $A^c$ is a finite disjoint union of sets in $\Sa$, i.e., $A^c \in \mathscr{B}$.

  We have shown that $\mathscr{B}$ is an algebra. Since $\Sa \sub \mathscr{B}$, $\A(\Sa) \sub \mathscr{B}$. Hence, if $A \in \A(\Sa)$, then $A \in \mathscr{B}$ and therefore $A$ can be expressed as a finite disjoint union of sets in $\Sa$.
\end{proof}